\section{Le cerveau comme source d'inspiration}

\begin{frame}{Le signal cérébral comme supervision}
\begin{itemize}
    \item Des enregistrements cérébraux servent à améliorer les modèles de langage
    \item Un modèle entraîné à résumer des récits devient plus aligné au cerveau
\end{itemize}
% TODO : vérifier et choisir 1 ou 2 articles (Toneva \& Wehbe 2019 ; Aw \& Toneva 2023)
\src{Toneva \& Wehbe, NeurIPS 2019 ; Aw \& Toneva, ICLR 2023}
\end{frame}

\begin{frame}{Échelle développementale : BabyLM}
\begin{itemize}
    \item Les modèles actuels sont entraînés sur des milliers de milliards de mots
    \item Le défi \textbf{BabyLM} limite l'entraînement à \textbf{10 à 100 millions de mots}
    \item Question : un modèle à échelle développementale est-il plus proche de l'humain ?
\end{itemize}
\src{Warstadt et al., BabyLM Challenge 2023}
\end{frame}
